\textbf{Proof.} Fix \(\beta\). We first compare the two adversarial classes. Substituting \(C=0\) in the historical-envelope restriction gives, for every pair of paths with \(z_t=z'_t\),
\[
\left|Y_t(z_{1:t})-Y_t(z'_{1:t})\right|
\le \sum_{\ell=1}^{t-1}0\cdot(1+\ell)^{-(1+\beta)}\mathbf 1\{z_{t-\ell}\ne z'_{t-\ell}\}=0.
\]
Hence \(Y_t(z_{1:t})=Y_t(z'_{1:t})\) whenever \(z_t=z'_t\). Together with \(\lvert Y_t(z_{1:t})\rvert\le B\), this is exactly the defining condition for \(\mathcal P_T^{(0)}\). Conversely, every \(Y\in\mathcal P_T^{(0)}\) has zero left-hand side for every same-current-treatment path pair and therefore satisfies the envelope with \(C=0\). Thus
\[
\mathcal P_{T,\beta}(0)=\mathcal P_T^{(0)}.
\]
Because the restriction compares only paths for which \(z_t=z'_t\), it imposes no equality between \(Y_t(1,z_{1:t-1})\) and \(Y_t(0,z'_{1:t-1})\); the contemporaneous treatment contrast remains unrestricted except for boundedness. The definitions of the two risks then have the same design class, estimator class, loss, target, and adversarial class, so
\[
R_T(\beta;0)=R_T^{(0)}.
\]

For the upper bound, assign \(W_1,\ldots,W_T\) independently as \(\operatorname{Bernoulli}(1/2)\) variables. This rule belongs to \(\mathcal D_T\). If \(Y\in\mathcal P_T^{(0)}\), let \(y_{t,a}\) be the common value of \(Y_t(z_{1:t})\) over paths with \(z_t=a\). Define the observed-data Horvitz--Thompson estimator
\[
\widehat\theta_{\mathrm{HT}}
=\frac1T\sum_{t=1}^T\left\{\frac{W_tY_t^{\mathrm{obs}}}{1/2}-\frac{(1-W_t)Y_t^{\mathrm{obs}}}{1/2}\right\}
=\frac1T\sum_{t=1}^T\{2W_ty_{t,1}-2(1-W_t)y_{t,0}\}.
\]
The expectation of its \(t\)-th unscaled summand is \(y_{t,1}-y_{t,0}\), and hence \(\widehat\theta_{\mathrm{HT}}\) is unbiased for \(\operatorname{GATE}_T(Y)\). Its centered \(t\)-th unscaled summand is
\[
(2W_t-1)(y_{t,1}+y_{t,0}).
\]
Independence across times and boundedness therefore imply
\[
\operatorname{Var}_\pi(\widehat\theta_{\mathrm{HT}})
=\frac1{T^2}\sum_{t=1}^T(y_{t,1}+y_{t,0})^2
\le \frac{4B^2}{T}.
\]
By Cauchy--Schwarz,
\[
\sup_{Y\in\mathcal P_T^{(0)}}\mathbb E_\pi\left|\widehat\theta_{\mathrm{HT}}-\operatorname{GATE}_T(Y)\right|
\le \frac{2B}{\sqrt T},
\]
so \(R_T^{(0)}\le 2B/\sqrt T\).

For the matching lower bound, fix an arbitrary \(\pi\in\mathcal D_T\) and an arbitrary observed-data estimator \(\widehat\theta\). As a Bayes proof device, draw mutually independent signs \(X_{t,a}\), for \(t=1,\ldots,T\) and \(a\in\{0,1\}\), each uniformly distributed on \(\{-B,B\}\), before assignment, and set
\[
Y_t(z_{1:t})=X_{t,z_t}.
\]
Every realized array is fixed throughout the experiment, is bounded by \(B\), and belongs to \(\mathcal P_T^{(0)}\). Thus the prior is supported on the finite-population adversarial class and introduces no outcome noise conditional on its realized array.

Let \(D=(W,Y^{\mathrm{obs}})\). Conditional on any realized \(D\) having positive probability, the unobserved coordinates \(X_{t,1-W_t}\), \(t=1,\ldots,T\), remain mutually independent and uniform on \(\{-B,B\}\). Indeed, the joint mass of the latent signs and \(D\) factors into the independent prior masses, indicators fixing the observed coordinates, and
\[
\prod_{t=1}^T\pi_t\bigl(W_t\mid W_{1:t-1},Y_{1:t-1}^{\mathrm{obs}}\bigr).
\]
The last factor is measurable with respect to \(D\) and contains no unobserved coordinate. This factorization applies equally to randomized, deterministic, history-dependent, and outcome-adaptive kernels and needs no positivity condition.

Define
\[
K(D)=\frac1T\sum_{t=1}^T(2W_t-1)Y_t^{\mathrm{obs}}.
\]
The conditional factorization gives
\[
\operatorname{GATE}_T(Y)\mid D
\ \overset{d}{=}\ K(D)+\frac BT S_T,
\qquad S_T=\sum_{t=1}^T\varepsilon_t,
\]
where, conditionally on \(D\), the variables \(\varepsilon_1,\ldots,\varepsilon_T\) are independent Rademacher signs. For every real \(a\), symmetry of \(S_T\) and the triangle inequality yield
\[
\mathbb E|a-S_T|
=\frac12\mathbb E\{|a-S_T|+|a+S_T|\}
\ge \mathbb E|S_T|.
\]
Applying this conditionally with \(a=T\{\widehat\theta-K(D)\}/B\) shows that
\[
\mathbb E\left[\left|\widehat\theta-\operatorname{GATE}_T(Y)\right|\mid D\right]
\ge \frac BT\mathbb E|S_T|.
\]

We now evaluate this Rademacher moment. If \(J\sim\operatorname{Binomial}(T,1/2)\), then \(S_T\) has the law of \(2J-T\). Put \(r=\lfloor T/2\rfloor+1\). Symmetry and the identity
\[
(2j-T)\binom Tj
=T\left\{\binom{T-1}{j-1}-\binom{T-1}{j}\right\}
\]
give the telescoping formula
\[
\mathbb E|S_T|
=2^{1-T}\sum_{j=r}^T(2j-T)\binom Tj
=T2^{-(T-1)}\binom{T-1}{\lfloor(T-1)/2\rfloor}.
\]
This quantity is at least \(\sqrt{T/2}\). To verify the central-binomial bound directly, set \(p_k=4^{-k}\binom{2k}{k}\) for \(k\ge1\). Since \(p_1=1/2\) and
\[
\frac{p_{k+1}}{p_k}=\frac{2k+1}{2k+2}\ge\sqrt{\frac{k}{k+1}},
\]
induction gives \(p_k\ge(2\sqrt{k})^{-1}\). Consequently, for every integer \(n\ge0\),
\[
2^{-n}\binom n{\lfloor n/2\rfloor}\ge\frac1{\sqrt{2(n+1)}};
\]
for even \(n=2k\) this follows from the bound on \(p_k\), for odd \(n=2k+1\) it follows from the bound on \(p_{k+1}\), and \(n=0\) is immediate. Taking \(n=T-1\) proves the claimed moment bound.

Averaging the conditional loss over \(D\), the Bayes risk under this supported prior is therefore at least \(B/\sqrt{2T}\). A supremum risk is no smaller than its average under any supported prior, so the arbitrary pair \((\pi,\widehat\theta)\) satisfies
\[
\sup_{Y\in\mathcal P_T^{(0)}}\mathbb E_\pi\left|\widehat\theta-\operatorname{GATE}_T(Y)\right|
\ge\frac{B}{\sqrt{2T}}.
\]
Taking the infima over \(\widehat\theta\) and \(\pi\) gives \(R_T^{(0)}\ge B/\sqrt{2T}\). Combining the two bounds,
\[
\frac{B}{\sqrt{2T}}
\le R_T(\beta;0)=R_T^{(0)}
\le\frac{2B}{\sqrt T}.
\]
Hence \(R_T(\beta;0)=R_T^{(0)}\asymp_B T^{-1/2}\), with the stated independent assignment rule and timewise Horvitz--Thompson estimator. \(\square\)